\section{Outcome (f): Conditional probability}

\subsection{What the outcome asks}

The exam tests whether you can read the phrase ``given that'' (or one of its
disguises) and respond by restricting the sample space to the conditioning event
before you count or divide. Questions supply the information in one of a few
forms: joint and marginal probabilities stated in words, a two-way table of
counts or percentages, a Venn setup with two or three coverages, or a small
population from which items are drawn. You compute $\Prob(A \mid B)$ from that
information, and you apply the complement and addition rules \emph{inside} the
conditional world when the question asks for ``not'' or ``or'' after ``given''.
Questions that give $\Prob(A\mid B)$ and ask for $\Prob(B\mid A)$ across several
classes are outcome (g) (Bayes and total probability); the definition that makes
those work is the subject of this section, and the two-class version is solved
here directly from the definition.

\subsection{Definition and properties}

\begin{keyfact}[Definition of conditional probability]
For events $A$ and $B$ with $\Prob(B) > 0$,
\[
\Prob(A \mid B) \;=\; \frac{\Prob(A \cap B)}{\Prob(B)} .
\]
Read: the probability of $A$ given that $B$ has occurred. The denominator is the
probability of the \emph{conditioning} event, the one named after ``given''.
\end{keyfact}

Interpretation. Conditioning on $B$ throws away every outcome outside $B$ and
treats $B$ as the new sample space. Probabilities of outcomes inside $B$ are then
rescaled by dividing by $\Prob(B)$ so that they again sum to $1$. Only the part
of $A$ that lies inside $B$, namely $A \cap B$, survives the restriction.

Because $\Prob(\,\cdot \mid B)$ is a genuine probability measure on the reduced
sample space $B$, it satisfies all three axioms: $\Prob(A\mid B) \ge 0$,
$\Prob(B \mid B) = 1$, and it is additive over disjoint events. Every rule
derived from the axioms therefore has a conditional version, obtained by
writing ``$\mid B$'' after every event.

\begin{keyfact}[Conditional complement and addition rules]
Holding the conditioning event $B$ fixed:
\begin{align*}
\Prob(\cc{A} \mid B) &= 1 - \Prob(A \mid B), \\
\Prob(A \cup C \mid B) &= \Prob(A \mid B) + \Prob(C \mid B) - \Prob(A \cap C \mid B), \\
\Prob(A \cap \cc{C} \mid B) &= \Prob(A \mid B) - \Prob(A \cap C \mid B).
\end{align*}
These hold because the event \emph{after} the bar is the same throughout.
\end{keyfact}

\begin{keyfact}[What is NOT true]
Changing the event \emph{after} the bar is not a complement operation:
\[
\Prob(A \mid \cc{B}) \;\ne\; 1 - \Prob(A \mid B) \quad \text{in general.}
\]
$\Prob(A\mid B)$ and $\Prob(A\mid \cc{B})$ live on two different reduced sample
spaces and need not be related at all. To get $\Prob(A \mid \cc{B})$ you need
$\Prob(A \cap \cc{B})$ and $\Prob(\cc{B})$, typically via
$\Prob(A \cap \cc{B}) = \Prob(A) - \Prob(A \cap B)$.
\end{keyfact}

\begin{keyfact}[Multiplication rule and chain rule]
Rearranging the definition gives
\[
\Prob(A \cap B) = \Prob(B)\,\Prob(A \mid B) = \Prob(A)\,\Prob(B \mid A),
\]
and for three events
\[
\Prob(A \cap B \cap C) = \Prob(A)\,\Prob(B \mid A)\,\Prob(C \mid A \cap B).
\]
Conversely, conditioning on an intersection is the same definition with a
compound conditioning event:
\[
\Prob(C \mid A \cap B) = \frac{\Prob(A \cap B \cap C)}{\Prob(A \cap B)} .
\]
\end{keyfact}

\begin{keyfact}[Link to independence]
If $\Prob(B) > 0$, then $A$ and $B$ are independent exactly when
$\Prob(A \mid B) = \Prob(A)$, equivalently $\Prob(A \mid B) = \Prob(A \mid \cc{B})$.
Knowing $B$ happened does not change the probability of $A$. If
$\Prob(A \mid B) \ne \Prob(A)$ the events are dependent, and the multiplication
rule must use the conditional probability, not the marginal.
\end{keyfact}

Two consequences. $\Prob(A \mid B)$ and $\Prob(B \mid A)$ share the numerator
$\Prob(A \cap B)$ but have different denominators, so
$\Prob(B \mid A) = \Prob(A \mid B)\,\Prob(B)/\Prob(A)$, which reverses a single
condition whenever $\Prob(A)$ is given. If $A \subseteq B$ then
$\Prob(A \mid B) = \Prob(A)/\Prob(B)$; if $A \cap B = \emptyset$ then
$\Prob(A \mid B) = 0$.

\subsection{Three computational settings}

\subsubsection*{(i) Two-way (contingency) table}

The conditioning event names a row or a column. Restrict to that row or column,
then divide the cell you want by the row or column total. Nothing outside that
row or column is used.

A table of counts. An insurer classifies $200$ life policyholders by sex and
smoking status.

\begin{center}
\begin{tabular}{@{}lrrr@{}}
\toprule
 & Smoker & Nonsmoker & Total \\
\midrule
Male   & 30 & 90  & 120 \\
Female & 24 & 56  & 80  \\
\midrule
Total  & 54 & 146 & 200 \\
\bottomrule
\end{tabular}
\end{center}

\begin{itemize}[nosep]
\item $\Prob(\text{Smoker} \mid \text{Male}) = 30/120 = 0.25$: restrict to the Male row.
\item $\Prob(\text{Smoker} \mid \text{Female}) = 24/80 = 0.30$: restrict to the Female row.
\item $\Prob(\text{Male} \mid \text{Smoker}) = 30/54 = 0.5556$: restrict to the Smoker column.
\item $\Prob(\text{Male} \cap \text{Smoker}) = 30/200 = 0.15$: no restriction; divide by the grand total.
\end{itemize}
Since $\Prob(\text{Smoker} \mid \text{Male}) \ne \Prob(\text{Smoker} \mid \text{Female})$,
sex and smoking status are dependent in this population.

A table of proportions. The same procedure works when the cells are joint
probabilities that sum to $1$; the row and column totals are then the marginal
probabilities.

\begin{center}
\begin{tabular}{@{}lrrr@{}}
\toprule
 & Claim & No claim & Total \\
\midrule
Young & 0.12 & 0.28 & 0.40 \\
Old   & 0.09 & 0.51 & 0.60 \\
\midrule
Total & 0.21 & 0.79 & 1.00 \\
\bottomrule
\end{tabular}
\end{center}

\begin{itemize}[nosep]
\item $\Prob(\text{Claim} \mid \text{Young}) = 0.12/0.40 = 0.30$.
\item $\Prob(\text{Claim} \mid \text{Old}) = 0.09/0.60 = 0.15$.
\item $\Prob(\text{Young} \mid \text{Claim}) = 0.12/0.21 = 0.5714$.
\item $\Prob(\text{No claim} \mid \text{Young}) = 1 - 0.30 = 0.70 = 0.28/0.40$: conditional complement, same row.
\item If the table had several age rows, conditioning on ``age at least 30''
would mean adding the relevant rows first: the denominator is the sum of those
row totals and the numerator is the sum of their Claim cells.
\end{itemize}

Exam problems often give the table indirectly: a marginal (``40\% of
policyholders are young''), a conditional (``30\% of young policyholders filed
a claim''), and an overall rate (``21\% filed a claim''). Rebuild the table:
Young--Claim $= 0.40 \times 0.30 = 0.12$, Old--Claim $= 0.21 - 0.12 = 0.09$,
the rest by subtraction. Every conditional is then a ratio of two entries.

\subsubsection*{(ii) Venn diagram / region bookkeeping}

With two or three overlapping coverages, label each disjoint region with its
probability (or count), then
\[
\Prob(A \mid B) = \frac{\text{total of the regions inside } A \cap B}{\text{total of the regions inside } B}.
\]
For three sets, the ``only'' regions come from inclusion--exclusion:
$\Prob(A \text{ only}) = \Prob(A) - \Prob(A\cap B) - \Prob(A \cap C) + \Prob(A \cap B \cap C)$.
``Of those with auto, what fraction also have home'' means the denominator is
every region inside the auto circle and the numerator is the lens $A \cap H$
(both of its regions, with and without life).

\subsubsection*{(iii) Sequential / tree}

When the experiment happens in stages (select a class, then observe a claim;
draw one file, then another), the branch labels on a probability tree are
conditional probabilities given the path so far. Multiplying along a path gives
the joint probability of the path. ``Given the first was $X$, what is the
probability the second is $Y$'' is read directly off a branch label; ``given the
second was $Y$'' reverses the tree, which is outcome (g).

\subsubsection*{Phrases that signal conditioning}

\begin{center}
\begin{tabular}{@{}p{0.42\textwidth}p{0.50\textwidth}@{}}
\toprule
Phrase in the problem & Meaning \\
\midrule
``given that $B$'' / ``if it is known that $B$'' & Denominator is $\Prob(B)$. \\
``of those who $B$'' / ``among $B$ policyholders'' & The $B$ group is the whole population; take the fraction inside it. \\
``restricted to $B$'' / ``within the $B$ class'' & Same. \\
``$x$\% of $B$ policyholders are $A$'' & $\Prob(A \mid B) = x/100$, not $\Prob(A \cap B)$. \\
``$x$\% of policyholders are both $A$ and $B$'' & Joint $\Prob(A \cap B)$; no conditioning. \\
``the probability that a $B$ policyholder is $A$'' & $\Prob(A \mid B)$. \\
\bottomrule
\end{tabular}
\end{center}

\subsection{Worked examples}

\begin{example}[Two-way table of policyholders]
An auto insurer has $500$ policyholders, classified by age and sex as follows.

\begin{center}
\begin{tabular}{@{}lrrr@{}}
\toprule
 & Male & Female & Total \\
\midrule
Young & 140 & 85  & 225 \\
Old   & 160 & 115 & 275 \\
\midrule
Total & 300 & 200 & 500 \\
\bottomrule
\end{tabular}
\end{center}

A policyholder is selected at random. Calculate the probability that the
policyholder is female, given that the policyholder is young.

\medskip
(A) 0.17 \quad (B) 0.38 \quad (C) 0.40 \quad (D) 0.43 \quad (E) 0.45
\end{example}

\begin{solution}
The conditioning event is ``young'', so restrict to the Young row, which has
$225$ policyholders. Of these, $85$ are female.
\[
\Prob(\text{Female} \mid \text{Young}) = \frac{85}{225} = 0.3778 .
\]
The distractors are the other ratios one could form from the same cell: $85/500 = 0.17$ is
$\Prob(\text{Female} \cap \text{Young})$; $200/500 = 0.40$ is the marginal $\Prob(\text{Female})$;
$85/200 = 0.425$ is the reversed conditional $\Prob(\text{Young} \mid \text{Female})$.
\textbf{Answer: (B).}
\end{solution}

\begin{example}[Both, given at least one]
For a randomly selected policyholder holding both auto and homeowners coverage
with an insurer, the probability of an auto claim in a given year is $0.30$, the
probability of a homeowners claim is $0.20$, and the probability of both is
$0.08$. Given that the policyholder files at least one claim in the year,
calculate the probability that the policyholder files both.

\medskip
(A) 0.08 \quad (B) 0.16 \quad (C) 0.19 \quad (D) 0.27 \quad (E) 0.42
\end{example}

\begin{solution}
Let $A$ = auto claim and $H$ = homeowners claim. ``At least one'' is $A \cup H$,
and ``both'' is $A \cap H$, which is a subset of $A \cup H$, so the intersection
of the two events is just $A \cap H$.
\[
\Prob(A \cup H) = 0.30 + 0.20 - 0.08 = 0.42, \qquad
\Prob(A \cap H \mid A \cup H) = \frac{0.08}{0.42} = 0.1905 .
\]
\textbf{Answer: (C).}
\end{solution}

\begin{example}[Venn with three coverages]
A survey of an insurer's customers finds that $55\%$ have auto coverage, $35\%$
have homeowners coverage, and $25\%$ have life coverage. Also, $15\%$ have both
auto and homeowners, $12\%$ have both auto and life, $8\%$ have both homeowners
and life, and $5\%$ have all three. Of those customers who have auto coverage,
calculate the proportion who have neither homeowners nor life coverage.

\medskip
(A) 0.33 \quad (B) 0.40 \quad (C) 0.45 \quad (D) 0.55 \quad (E) 0.60
\end{example}

\begin{solution}
Let $A$, $H$, $L$ be the three coverages. The denominator is $\Prob(A) = 0.55$.
The numerator is the ``auto only'' region:
\[
\Prob(A \cap \cc{H} \cap \cc{L}) = \Prob(A) - \Prob(A\cap H) - \Prob(A \cap L) + \Prob(A\cap H \cap L)
= 0.55 - 0.15 - 0.12 + 0.05 = 0.33 .
\]
Then
\[
\Prob(\cc{H} \cap \cc{L} \mid A) = \frac{0.33}{0.55} = 0.60 .
\]
Choice (A) is the joint probability $0.33$, not the conditional.
\textbf{Answer: (E).}
\end{solution}

\begin{example}[Reversing one condition from the definition]
In a block of life policies, $30\%$ of policyholders are smokers. Among
smokers, $20\%$ filed a claim last year. Overall, $12\%$ of policyholders filed
a claim last year. Calculate the probability that a policyholder who filed a
claim last year is a smoker.

\medskip
(A) 0.06 \quad (B) 0.20 \quad (C) 0.30 \quad (D) 0.50 \quad (E) 0.60
\end{example}

\begin{solution}
Let $S$ = smoker and $C$ = filed a claim. Given: $\Prob(S) = 0.30$,
$\Prob(C \mid S) = 0.20$, $\Prob(C) = 0.12$. The joint probability is
\[
\Prob(S \cap C) = \Prob(S)\,\Prob(C \mid S) = 0.30 \times 0.20 = 0.06 ,
\]
and the requested conditional is
\[
\Prob(S \mid C) = \frac{\Prob(S \cap C)}{\Prob(C)} = \frac{0.06}{0.12} = 0.50 .
\]
No Bayes machinery is needed because $\Prob(C)$ is supplied directly; the
definition does the whole job.
\textbf{Answer: (D).}
\end{solution}

\begin{example}[Conditioning on the complement]
A life insurer classifies $60\%$ of its policyholders as standard and the rest
as preferred. The probability that a standard policyholder lapses within one
year is $0.15$. The probability that a randomly selected policyholder lapses
within one year is $0.11$. Calculate the probability that a preferred
policyholder lapses within one year.

\medskip
(A) 0.02 \quad (B) 0.05 \quad (C) 0.085 \quad (D) 0.15 \quad (E) 0.85
\end{example}

\begin{solution}
Let $S$ = standard and $L$ = lapse. Given: $\Prob(S) = 0.60$, so
$\Prob(\cc{S}) = 0.40$; $\Prob(L \mid S) = 0.15$; $\Prob(L) = 0.11$.
Build the pieces of $L$:
\[
\Prob(L \cap S) = 0.60 \times 0.15 = 0.09, \qquad
\Prob(L \cap \cc{S}) = \Prob(L) - \Prob(L \cap S) = 0.11 - 0.09 = 0.02 .
\]
Then
\[
\Prob(L \mid \cc{S}) = \frac{0.02}{0.40} = 0.05 .
\]
Choice (E), $0.85 = 1 - 0.15$, is $\Prob(\cc{L} \mid S)$, the probability that
a \emph{standard} policyholder does \emph{not} lapse; it says nothing about
preferred policyholders. Choice (A) is the joint probability, not the
conditional.
\textbf{Answer: (B).}
\end{solution}

\begin{example}[Conditioning on an intersection]
An insurer's records show that $50\%$ of its customers have auto coverage,
$40\%$ have homeowners coverage, and $20\%$ have both. Also, $12\%$ of all
customers filed a claim last year, $8\%$ have auto coverage and filed a claim
last year, and $5\%$ have both coverages and filed a claim last year. Calculate
the probability that a customer who has both auto and homeowners coverage filed
a claim last year.

\medskip
(A) 0.05 \quad (B) 0.10 \quad (C) 0.16 \quad (D) 0.25 \quad (E) 0.40
\end{example}

\begin{solution}
Let $A$ = auto, $H$ = homeowners, $C$ = claim. The conditioning event is
$A \cap H$ with $\Prob(A \cap H) = 0.20$; the numerator is $\Prob(A \cap H \cap C) = 0.05$.
\[
\Prob(C \mid A \cap H) = \frac{0.05}{0.20} = 0.25 .
\]
The other figures are distractors: $\Prob(C \mid A) = 0.08/0.50 = 0.16$ conditions
on auto alone; $\Prob(C) = 0.12$ and $\Prob(H) = 0.40$ are not needed.
\textbf{Answer: (D).}
\end{solution}

\subsection{Traps}

\begin{trap}[Dividing by the wrong marginal]
The denominator is always the probability of the event named after ``given''
(or ``of those who'', ``among''). In a table, that means the total of the row
or column you restricted to, never the grand total and never the other
margin. Before dividing, say out loud which group is the new population.
\end{trap}

\begin{trap}[$\Prob(A \mid B)$ versus $\Prob(B \mid A)$]
These share the numerator $\Prob(A \cap B)$ but have different denominators
and are generally different numbers. ``$20\%$ of smokers filed a claim'' is
$\Prob(C \mid S)$; ``$50\%$ of claimants are smokers'' is $\Prob(S \mid C)$.
Exam answer choices routinely include the reversed conditional as a
distractor. Write the conditional with the bar before computing anything.
\end{trap}

\begin{trap}[$\Prob(A \mid \cc{B})$ is not $1 - \Prob(A \mid B)$]
The complement rule flips the event \emph{before} the bar:
$\Prob(\cc{A} \mid B) = 1 - \Prob(A \mid B)$. Flipping the event \emph{after} the
bar changes the population, and the two conditionals need not add to $1$.
To find $\Prob(A \mid \cc{B})$, compute $\Prob(A \cap \cc{B}) = \Prob(A) - \Prob(A \cap B)$
and divide by $\Prob(\cc{B}) = 1 - \Prob(B)$.
\end{trap}

\begin{trap}[Joint versus conditional in the wording]
``$x$\% of policyholders are young and filed a claim'' is a joint probability
$\Prob(Y \cap C)$. ``$x$\% of young policyholders filed a claim'' is a
conditional $\Prob(C \mid Y)$. ``$x$\% of policyholders are young'' is a
marginal $\Prob(Y)$. Misreading one of these as another produces an answer
that is present among the choices. When the problem mixes the three, rebuild the
two-way table from them before answering.
\end{trap}

\begin{trap}[Mixing counts and percentages]
A conditional probability is a ratio, so it can be computed from counts or from
proportions, but not from a mixture. If a row total is a count and a cell is a
percentage of the whole population, convert both to the same units first. In
a proportion table, the grand total is $1$, and the marginals are row and
column sums; do not divide a proportion by a count.
\end{trap}

\subsection{Practice problems}

\begin{practice}
A life insurer classifies $400$ policyholders as preferred or standard, and as
smokers or nonsmokers. There are $20$ preferred smokers, $140$ preferred
nonsmokers, $90$ standard smokers, and $150$ standard nonsmokers. A
policyholder is selected at random. Calculate the probability that the
policyholder is preferred, given that the policyholder is a nonsmoker.

\medskip
(A) 0.35 \quad (B) 0.40 \quad (C) 0.48 \quad (D) 0.72 \quad (E) 0.88
\end{practice}

\begin{practice}
For a randomly selected employee with a group benefits plan, the probability of
a dental claim in a year is $0.45$, the probability of a vision claim is
$0.35$, and the probability of both is $0.15$. Given that the employee has at
least one of the two types of claim in the year, calculate the probability that
the employee has both.

\medskip
(A) 0.15 \quad (B) 0.23 \quad (C) 0.33 \quad (D) 0.43 \quad (E) 0.65
\end{practice}

\begin{practice}
In a survey of an insurer's customers, $40\%$ have auto coverage, $30\%$ have
homeowners coverage, and $25\%$ have life coverage; $12\%$ have auto and
homeowners, $10\%$ have auto and life, $8\%$ have homeowners and life, and
$4\%$ have all three. Of those customers with homeowners coverage, calculate the
proportion who have neither auto nor life coverage.

\medskip
(A) 0.14 \quad (B) 0.40 \quad (C) 0.47 \quad (D) 0.53 \quad (E) 0.60
\end{practice}

\begin{practice}
An auto insurer classifies $20\%$ of its policies as high risk. The probability
that a high-risk policy has a claim in a year is $0.35$. The probability that a
randomly selected policy has a claim in a year is $0.16$. Calculate the
probability that a policy with a claim in the year is a high-risk policy.

\medskip
(A) 0.07 \quad (B) 0.20 \quad (C) 0.35 \quad (D) 0.44 \quad (E) 0.56
\end{practice}

\begin{practice}
Of an insurer's policyholders, $55\%$ are female. The probability that a
female policyholder files a claim in a year is $0.10$. The probability that a
randomly selected policyholder files a claim in a year is $0.145$. Calculate
the probability that a male policyholder files a claim in the year.

\medskip
(A) 0.09 \quad (B) 0.145 \quad (C) 0.20 \quad (D) 0.30 \quad (E) 0.90
\end{practice}

\begin{practice}
An insurer's records show that $60\%$ of customers have auto coverage. Among
customers with auto coverage, $50\%$ carry collision coverage. Also, $12\%$ of
all customers have auto coverage with collision and filed a claim last year.
Calculate the probability that a customer with auto coverage and collision
coverage filed a claim last year.

\medskip
(A) 0.12 \quad (B) 0.20 \quad (C) 0.30 \quad (D) 0.40 \quad (E) 0.72
\end{practice}

\begin{practice}
A file cabinet contains $12$ claim files, of which $5$ are fraudulent. An
auditor selects $2$ files at random without replacement. Given that at least
one of the selected files is fraudulent, calculate the probability that both
are fraudulent.

\medskip
(A) 0.15 \quad (B) 0.22 \quad (C) 0.36 \quad (D) 0.41 \quad (E) 0.68
\end{practice}

\subsection{Answers to practice problems}

\begin{enumerate}[label=\arabic*., nosep]
\item \textbf{(C).} Nonsmokers total $140 + 150 = 290$; of these $140$ are preferred.
$140/290 = 0.4828$. ($140/400 = 0.35$ is the joint; $160/400 = 0.40$ is the marginal;
$140/160 = 0.875$ is the reversed conditional.)
\item \textbf{(B).} $\Prob(D \cup V) = 0.45 + 0.35 - 0.15 = 0.65$;
$\Prob(D \cap V \mid D \cup V) = 0.15/0.65 = 0.2308$.
\item \textbf{(C).} Homeowners only $= 0.30 - 0.12 - 0.08 + 0.04 = 0.14$;
divide by $\Prob(H) = 0.30$: $0.14/0.30 = 0.4667$. ($0.12/0.30 = 0.40$ is the
fraction of homeowners customers with auto.)
\item \textbf{(D).} $\Prob(R \cap C) = 0.20 \times 0.35 = 0.07$;
$\Prob(R \mid C) = 0.07/0.16 = 0.4375$.
\item \textbf{(C).} $\Prob(F \cap C) = 0.55 \times 0.10 = 0.055$;
$\Prob(M \cap C) = 0.145 - 0.055 = 0.09$; $\Prob(C \mid M) = 0.09/0.45 = 0.20$.
(Choice (E), $0.90 = 1 - 0.10$, is $\Prob(\cc{C} \mid F)$, not a statement about males.)
\item \textbf{(D).} $\Prob(A \cap K) = 0.60 \times 0.50 = 0.30$;
$\Prob(C \mid A \cap K) = 0.12/0.30 = 0.40$.
\item \textbf{(B).} $\Prob(\text{both}) = \binom{5}{2}/\binom{12}{2} = 10/66$;
$\Prob(\text{none}) = \binom{7}{2}/\binom{12}{2} = 21/66$;
$\Prob(\text{at least one}) = 45/66$;
ratio $= 10/45 = 2/9 = 0.2222$. ($4/11 = 0.36$ is $\Prob(\text{second fraudulent} \mid \text{first fraudulent})$, a different condition.)
\end{enumerate}

\subsection{One-page summary}

\begin{itemize}[nosep]
\item Definition: $\Prob(A \mid B) = \Prob(A \cap B)/\Prob(B)$, $\Prob(B) > 0$. The
event after the bar is the new sample space; divide by its probability.
\item Signal words for conditioning: given, of those who, among, if it is known
that, restricted to, within the class, ``$x$\% of $B$ policyholders are $A$''.
\item ``$x$\% of policyholders are $A$ and $B$'' is joint; ``$x$\% of $B$
policyholders are $A$'' is conditional; ``$x$\% of policyholders are $B$'' is marginal.
\item Two-way table: restrict to the row or column named after ``given'', divide
the cell by that row or column total. Works for counts or proportions, not a mix.
\item Rebuild the table from words: cell $= \text{marginal} \times \text{conditional}$;
remaining cells by subtraction from marginals.
\item Venn: $\Prob(A \mid B) = (\text{regions in } A \cap B)/(\text{regions in } B)$;
``only'' regions by inclusion--exclusion.
\item Tree: branch labels are conditionals; multiply along a path for the joint.
\item Conditional complement: $\Prob(\cc{A} \mid B) = 1 - \Prob(A \mid B)$.
\item Conditional addition: $\Prob(A \cup C \mid B) = \Prob(A \mid B) + \Prob(C \mid B) - \Prob(A \cap C \mid B)$.
\item Not a rule: $\Prob(A \mid \cc{B}) \ne 1 - \Prob(A \mid B)$. Get it from
$\Prob(A \cap \cc{B}) = \Prob(A) - \Prob(A \cap B)$ divided by $1 - \Prob(B)$.
\item Multiplication: $\Prob(A \cap B) = \Prob(B)\Prob(A \mid B) = \Prob(A)\Prob(B \mid A)$;
chain: $\Prob(A \cap B \cap C) = \Prob(A)\Prob(B \mid A)\Prob(C \mid A \cap B)$.
\item Conditioning on an intersection: $\Prob(C \mid A \cap B) = \Prob(A \cap B \cap C)/\Prob(A \cap B)$.
\item Reversing one condition: $\Prob(B \mid A) = \Prob(A \mid B)\Prob(B)/\Prob(A)$ when $\Prob(A)$ is given.
\item Both given at least one: $\Prob(A \cap B)/\Prob(A \cup B)$ with
$\Prob(A \cup B) = \Prob(A) + \Prob(B) - \Prob(A \cap B)$.
\item Independence: $\Prob(A \mid B) = \Prob(A) = \Prob(A \mid \cc{B})$.
\item Subset: $A \subseteq B \Rightarrow \Prob(A \mid B) = \Prob(A)/\Prob(B)$;
disjoint: $\Prob(A \mid B) = 0$.
\end{itemize}
